\section{AHSME 1983}
        \begin{enumerate}
        	\item (\textit{AHSME 1983}) If $x \neq 0, \frac x{2} = y^2$ and $\frac{x}{4} = 4y$, then $x$ equals


 $\textbf{(A)}\ 8\qquad \textbf{(B)}\ 16\qquad \textbf{(C)}\ 32\qquad \textbf{(D)}\ 64\qquad \textbf{(E)}\ 128$


 
 \textbf{Answer: }E



\textbf{Solution}

From $\frac{x}{4} = 4y$, we get $x=16y$. Plugging in the other equation, $\frac{16y}{2} = y^2$, so $y^2-8y=0$. Factoring, we get $y(y-8)=0$, so the solutions are $0$ and $8$. Since $x \neq 0$, we also have $y \neq 0$, so $y=8$. Hence $x=16\cdot 8 = \boxed{\textbf{(E)}\ 128}$.

 


	\item (\textit{AHSME 1983}) Point $P$ is outside circle $C$ on the plane. At most how many points on $C$ are $3$ cm from $P$? 


 $\textbf{(A)} \ 1 \qquad  \textbf{(B)} \ 2 \qquad  \textbf{(C)} \ 3 \qquad  \textbf{(D)} \ 4 \qquad  \textbf{(E)} \ 8$


 
 \textbf{Answer: }B



\textbf{Solution}

The points $3$ cm away from $P$ can be represented as a circle centered at $P$ with radius $3$ cm. The maximum number of intersection points of two circles is $\boxed{\textbf{(B)} \ 2}$.

 


	\item (\textit{AHSME 1983}) Three primes $p,q$ and $r$ satisfy $p+q = r$ and $1 < p < q$. Then $p$ equals


 $\textbf{(A)}\ 2\qquad \textbf{(B)}\ 3\qquad \textbf{(C)}\ 7\qquad \textbf{(D)}\ 13\qquad \textbf{(E)}\ 17$


 
 \textbf{Answer: }A



\textbf{Solution}

We are given that $p,q$ and $r$ are primes. In order for $p$ and $q$ to sum to another prime, either $p$ or $q$ has to be even, because the sum of two odd numbers would be even, and the only even prime is $2$ (but $p + q = 2$ would have, as the only solution in positive integers, $p = q = 1$, and $1$ is not prime). Thus, with one of either $p$ or $q$ being even, either $p$ or $q$ must be $2$, and as $p < q$, we deduce $p = 2$ (as $2$ is the smallest prime). This means the answer is $\boxed{\textbf{(A)}\ 2}$.

 


	\item (\textit{AHSME 1983}) 
                    \begin{figure}[H]
                        centering
                        \includegraphics[width=0.5\textwidth]{Problem/AHSME_I_1983_4_Problem_Figure_1.png}
                    \end{figure}
                    


 In the adjoining plane figure, sides $AF$ and $CD$ are parallel, as are sides $AB$ and $EF$, 
and sides $BC$ and $ED$. Each side has length $1$. Also, $\angle FAB = \angle BCD = 60^\circ$. 
The area of the figure is 


 $\textbf{(A)} \ \frac{\sqrt 3}{2} \qquad  \textbf{(B)} \ 1 \qquad  \textbf{(C)} \ \frac{3}{2} \qquad  \textbf{(D)}\ \sqrt{3}\qquad \textbf{(E)}\ 2$


 
 \textbf{Answer: }D



\textbf{Solution}


                    \begin{figure}[H]
                        centering
                        \includegraphics[width=0.5\textwidth]{Solution/AHSME_I_1983_4_Solution_Figure_1.png}
                    \end{figure}
                    

 By rotating the diagram and drawing the dotted lines, we see that the figure can be divided into four equilateral triangles, each of side length $1$. The area of one such equilateral triangle is $\frac{\sqrt{3}}{4} \cdot 1^2 = \frac{\sqrt{3}}{4}$, which gives a total of $4\left(\frac{\sqrt{3}}{4}\right) = \sqrt{3}$, or $\boxed{D}$.

 


	\item (\textit{AHSME 1983}) Triangle $ABC$ has a right angle at $C$. If $\sin A = \frac{2}{3}$, then $\tan B$ is


 $\textbf{(A)}\ \frac{3}{5}\qquad \textbf{(B)}\ \frac{\sqrt 5}{3}\qquad \textbf{(C)}\ \frac{2}{\sqrt 5}\qquad \textbf{(D)}\ \frac{\sqrt{5}}{2}\qquad \textbf{(E)}\ \frac{5}{3}$


 
 \textbf{Answer: }D



\textbf{Solution}

Since $\sin$ can be seen as "opposite over hypotenuse" in a right triangle, we can label the diagram as shown.
                    \begin{figure}[H]
                        centering
                        \includegraphics[width=0.5\textwidth]{Solution/AHSME_I_1983_5_Solution_Figure_1.png}
                    \end{figure}
                    By the Pythagorean Theorem, we have:
 \[y^2+(2x)^2=(3x)^2\]
 \[y=\sqrt{9x^2-4x^2}\]
 \[y=\sqrt{5x^2}\]
 \[y=x\sqrt{5}\]so $\tan{B} = \frac{x \sqrt{5}}{2x} = \frac{\sqrt{5}}{2}$, which is choice $\boxed{\textbf{(D)}}$.

 


	\item (\textit{AHSME 1983}) When $x^5, x+\frac{1}{x}$ and $1+\frac{2}{x} + \frac{3}{x^2}$ are multiplied, the product is a polynomial of degree


 $\textbf{(A)}\ 2\qquad \textbf{(B)}\ 3\qquad \textbf{(C)}\ 6\qquad \textbf{(D)}\ 7\qquad \textbf{(E)}\ 8$


 
 \textbf{Answer: }C



\textbf{Solution}

We have $x^5\left(x+\frac{1}{x}\right)\left(1+\frac{2}{x}+\frac{3}{x^2}\right) = (x^6+x^4)\left(1+\frac{2}{x}+\frac{3}{x^2}\right) = x^6 + \text{lower order terms}$, where we know that the $x^6$ will not get cancelled out by e.g. a $-x^6$ term since all the terms inside the brackets are positive. Thus the degree is $6$, which is choice $\boxed{\textbf{(C)}}$.

 


	\item (\textit{AHSME 1983}) Alice sells an item at $\$10$ less than the list price and receives $10\%$ of her selling price as her commission. 
Bob sells the same item at $\$20$ less than the list price and receives $20\%$ of his selling price as his commission. 
If they both get the same commission, then the list price is


 $\textbf{(A) } \$20\qquad \textbf{(B) } \$30\qquad \textbf{(C) } \$50\qquad \textbf{(D) } \$70\qquad \textbf{(E) } \$100$


 
 \textbf{Answer: }B



\textbf{Solution}

If $x$ is the list price, then $0.1(x-10)=0.2(x-20)$. Solving this equation gives $x=30$, so the answer is $\boxed{\textbf{(B) }\$30}$.

 


	\item (\textit{AHSME 1983}) Let $f(x) = \frac{x+1}{x-1}$. Then for $x^2 \neq 1, f(-x)$ is


 $\textbf{(A)}\ \frac{1}{f(x)}\qquad \textbf{(B)}\ -f(x)\qquad \textbf{(C)}\ \frac{1}{f(-x)}\qquad \textbf{(D)}\ -f(-x)\qquad \textbf{(E)}\ f(x)$


 
 \textbf{Answer: }A



\textbf{Solution}

We find $f(-x) = \frac{-x+1}{-x-1} = \frac{x-1}{x+1} = \frac{1}{f(x)}$, so the answer is $\boxed{\textbf{(A)}}$.

 


	\item (\textit{AHSME 1983}) In a certain population the ratio of the number of women to the number of men is $11$ to $10$. 
If the average (arithmetic mean) age of the women is $34$ and the average age of the men is $32$, 
then the average age of the population is


 $\textbf{(A)}\ 32\frac{9}{10}\qquad \textbf{(B)}\ 32\frac{20}{21}\qquad \textbf{(C)}\ 33\qquad \textbf{(D)}\ 33\frac{1}{21}\qquad \textbf{(E)}\ 33\frac{1}{10}$


 
 \textbf{Answer: }D



\textbf{Solution}

Assume, without loss of generality, that there are exactly $11$ women and $10$ men. Then the total age of the women is $34 \cdot 11 = 374$ and the total age of the men is $32 \cdot 10 = 320$. Therefore the overall average is $\frac{374+320}{11+10} = \boxed{\textbf{(D)}\ 33\frac{1}{21}}$.

 


	\item (\textit{AHSME 1983}) Segment $AB$ is both a diameter of a circle of radius $1$ and a side of an equilateral triangle $ABC$. 
The circle also intersects $AC$ and $BC$ at points $D$ and $E$, respectively. The length of $AE$ is 


 $\textbf{(A)} \ \frac{3}{2} \qquad  \textbf{(B)} \ \frac{5}{3} \qquad  \textbf{(C)} \ \frac{\sqrt 3}{2} \qquad  \textbf{(D)}\ \sqrt{3}\qquad \textbf{(E)}\ \frac{2+\sqrt 3}{2}$


 
 \textbf{Answer: }D



\textbf{Solution}

Note that since $AB$ is a diameter, $\angle AEB = 90^{\circ}$, which means $AB$ is an altitude of equilateral triangle $ABC$. It follows that $\triangle ABE$ is a $30^{\circ}-60^{\circ}-90^{\circ}$ triangle, and so $AE = AB \cdot \frac{\sqrt{3}}{2} = (2 \cdot 1) (\frac{\sqrt{3}}{2}) = \boxed{\textbf{(D)}\ \sqrt{3}}$.

 


	\item (\textit{AHSME 1983}) Simplify $\sin (x-y) \cos y + \cos (x-y) \sin y$.


 $\textbf{(A)}\ 1\qquad \textbf{(B)}\ \sin x\qquad \textbf{(C)}\ \cos x\qquad \textbf{(D)}\ \sin x \cos 2y\qquad \textbf{(E)}\ \cos x\cos 2y$


 
 \textbf{Answer: }B



\textbf{Solution}

By the addition formula for $\sin$, this becomes $\sin{((x-y)+y)} = \sin{x}$, so the answer is $\boxed{\textbf{(B)}}$.

 


	\item (\textit{AHSME 1983}) If $\log_7 \Big(\log_3 (\log_2 x) \Big) = 0$, then $x^{-1/2}$ equals 


 $\textbf{(A)} \ \frac{1}{3} \qquad  \textbf{(B)} \ \frac{1}{2 \sqrt 3} \qquad  \textbf{(C)}\ \frac{1}{3\sqrt 3}\qquad \textbf{(D)}\ \frac{1}{\sqrt{42}}\qquad \textbf{(E)}\ \text{none of these}$


 
 \textbf{Answer: }E



\textbf{Solution}

Because $\log_7 \Big(\log_3 (\log_2 x) \Big) = 0$, we deduce $\log_3 (\log_2 x) =1$, and thus $\log_2 x=3$. Therefore, $x=8$, which means $x^{-1/2}=\frac{1}{2\sqrt{2}}$. Since this does not match any of the answer choices, the answer is $\fbox{{\bf(E)} \text{none of these}}$.

 


	\item (\textit{AHSME 1983}) If $xy = a, xz =b,$ and $yz = c$, and none of these quantities is $0$, then $x^2+y^2+z^2$ equals


 $\textbf{(A)}\ \frac{ab+ac+bc}{abc}\qquad \textbf{(B)}\ \frac{a^2+b^2+c^2}{abc}\qquad \textbf{(C)}\ \frac{(a+b+c)^2}{abc}\qquad \textbf{(D)}\ \frac{(ab+ac+bc)^2}{abc}\qquad \textbf{(E)}\ \frac{(ab)^2+(ac)^2+(bc)^2}{abc}$


 
 \textbf{Answer: }E



\textbf{Solution}

From the equations, we deduce $x = \frac{a}{y}, z = \frac{b}{x},$ and $y = \frac{c}{z}$. Substituting these into the expression $x^2+y^2+z^2$ thus gives $\frac{a^2}{y^2} + \frac{b^2}{x^2} + \frac{c^2}{z^2} = \frac{a^2x^2z^2+b^2y^2z^2+c^2y^2x^2}{x^2y^2z^2} = \frac{a^2b^2+b^2c^2+c^2a^2}{x^2y^2z^2}$, so the answer is $\boxed{\textbf{(E)}\ \frac{(ab)^2+(ac)^2+(bc)^2}{abc}}$.

 



\textbf{Solution 2}

$x^2$ is $\frac{abc}{c^2}$, $y^2$ is $\frac{abc}{b^2}$, and $z^2$ is $\frac{abc}{a^2}$, so $x^2+y^2+z^2$ is $\frac{abc}{c^2}+\frac{abc}{b^2}+\frac{abc}{a^2}= \boxed{\textbf{(E)}\ \frac{(ab)^2+(ac)^2+(bc)^2}{abc}}$

 -purplepenguin2

 


	\item (\textit{AHSME 1983}) The units digit of $3^{1001} 7^{1002} 13^{1003}$ is


 $\textbf{(A)}\ 1\qquad \textbf{(B)}\ 3\qquad \textbf{(C)}\ 5\qquad \textbf{(D)}\ 7\qquad \textbf{(E)}\ 9$


 
 \textbf{Answer: }E



\textbf{Solution 1}

First, we notice that $3^0$ is congruent to $1 \ \text{(mod 10)}$, $3^1$ is $3 \ \text{(mod 10)}$, $3^2$ is $9 \ \text{(mod 10)}$, $3^3$ is $7 \ \text{(mod 10)}$, $3^4$ is $1 \ \text{(mod 10)}$, and so on. This turns out to be a cycle repeating every $4$ terms, so $3^{1001}$ is congruent to $3 \ \text{(mod 10)}$.

 The number $7$ has a similar cycle, going $1, 7, 9, 3, 1, ...$ Hence we have that $7^{1002}$ is congruent to $9 \ \text{(mod 10)}$. Finally, $13^{1003}$ is congruent to $3^{1003} \equiv 7 \ \text{(mod 10)}$. Thus the required units digit is $3\cdot 9\cdot 7 \equiv 9 \ \text{(mod 10)}$, so the answer is $\boxed{\textbf{(E)}\ 9}$.

 



\textbf{Solution 2}

By Euler's Totient Theorem, if $\gcd(a, 10) = 1$, then $a^{\phi(10)} \equiv a^4 \equiv 1\ (\mathrm{mod}\ 10)$.  So $3^{1001} \cdot 7^{1002} \cdot 13^{1003} \equiv 3^{1001} \cdot 7^{1002} \cdot 3^{1003} \equiv 3^{2004} \cdot 7^{1002} \equiv 3^0 \cdot 7^2 \equiv 1 \cdot 9 \equiv \boxed{\textbf{(E)}\ 9}\ (\mathrm{mod}\ 10)$.

 -j314andrews

 


	\item (\textit{AHSME 1983}) Three balls marked $1,2$ and $3$ are placed in an urn. One ball is drawn, its number is recorded, and then the ball is returned to the urn. This process is repeated and then repeated once more, and each ball is equally likely to be drawn on each occasion. If the sum of the numbers recorded is $6$, what is the probability that the ball numbered $2$ was drawn all three times? 


 $\textbf{(A)} \ \frac{1}{27} \qquad  \textbf{(B)} \ \frac{1}{8} \qquad  \textbf{(C)} \ \frac{1}{7} \qquad  \textbf{(D)} \ \frac{1}{6} \qquad \textbf{(E)}\ \frac{1}{3}$


 
 \textbf{Answer: }C



\textbf{Solution}

Since a ball is drawn three times and the sum is $6$, the only ways this can happen are the permutations of $1, 2, 3$ and $2, 2, 2$. Therefore, with there being $3! + 1 = 7$ equally-likely possibilities in total, the probability that the ball numbered $2$ was drawn all three times is $\boxed{\textbf{(C)}\ \frac{1}{7}}$.

 


	\item (\textit{AHSME 1983}) Let $x = .123456789101112....998999$, where the digits are obtained by writing the integers $1$ through $999$ in order. 
The $1983$<sup>rd</sup> digit to the right of the decimal point is


 $\textbf{(A)}\ 2\qquad \textbf{(B)}\ 3\qquad \textbf{(C)}\ 5\qquad \textbf{(D)}\ 7\qquad \textbf{(E)}\ 8$


 
 \textbf{Answer: }D



\textbf{Solution}

We consider the first $1983$ digits, letting the $1983$<sup>rd</sup> digit be $z$. We can break the string of digits into three segments: let $A$ denote $123456789$ (the $1$-digit numbers), let $B$ denote $1011...9899$ (the $2$-digit numbers), and let $C$ denote $100101...z$ (the $3$-digit numbers). Clearly there are $9$ digits in $A$; in $B$, there are $99-10+1 = 90$ numbers, so $90 \cdot 2 = 180$ digits. This leaves $1983 - 9 - 180 = 1794$ digits in $C$. Notice that $1794 = 3 \cdot 598$ with no remainder, so $C$ consists of precisely the first $598$ $3$-digit numbers. Since the first $3$-digit number is $100$, the $598$<sup>th</sup> is $100 + 598 - 1 = 697$, so as $z$ is the last digit, the answer is $\boxed{\textbf{(D)}\ 7}$.

 


	\item (\textit{AHSME 1983}) 
                    \begin{figure}[H]
                        centering
                        \includegraphics[width=0.5\textwidth]{Problem/AHSME_I_1983_17_Problem_Figure_1.png}
                    \end{figure}
                    


 The diagram above shows several numbers in the complex plane. The circle is the unit circle centered at the origin. 
One of these numbers is the reciprocal of $F$. Which one? 


 $\textbf{(A)} \ A \qquad  \textbf{(B)} \ B \qquad  \textbf{(C)} \ C \qquad  \textbf{(D)} \ D \qquad  \textbf{(E)} \ E$


 
 \textbf{Answer: }C



\textbf{Solution 1 (Linear Modeling)}

The reciprocal of a number \( a \) is \( \frac{1}{a} \).

 We start with \( F \) and the point \( (1, 0i) \) where \( i = \sqrt{-1} \). Draw a line through these two points. 

 The reciprocal of a line is legitimately a vertical shrink or a horizontal stretch (the line gets farther away from the \( i_1 \) axis but closer to the real axis). Therefore the line must pass through point $\boxed{\textbf{C}}$ as it is the only point that is closest to the real axis whilst maintaining a reflection over the imaginary axis.

 



\textbf{Solution 2}

Write $F$ as $a + bi$, where we see from the diagram that $a, b > 0$ and $a^2+b^2>1$ (as $F$ is outside the unit circle). We have $\frac{1}{a+bi} = \frac{a-bi}{a^2+b^2} = \frac{a}{a^2+b^2} - \frac{b}{a^2+b^2}i$, so, since $a, b > 0$, the reciprocal of $F$ has a positive real part and negative imaginary part. Also, the reciprocal has magnitude equal to the reciprocal of $F$'s magnitude (since $|a||b| = |ab|$); as $F$'s magnitude is greater than $1$, its reciprocal's magnitude will thus be between $0$ and $1$, so its reciprocal will be inside the unit circle. Therefore, the only point shown which could be the reciprocal of $F$ is point $\boxed{\textbf{C}}$.

 



\textbf{Solution 3 (Polar Form)}

Let $z = r(\cos\theta + i\sin\theta)$ be the complex number at $F$ in polar form.  Then $\frac{1}{z} = \frac{1}{r}(\cos(-\theta) + i\sin(-\theta))$. Since $F$ is outside the circle, $r > 1$, and therefore $\frac{1}{r} < 1$.  So $\frac{1}{z}$ must be inside the circle and on the line which is $\theta$ clockwise from the real axis, that is $\boxed{(\mathbf{C})\ C}$.

 -j314andrews

 


	\item (\textit{AHSME 1983}) Let $f$ be a polynomial function such that, for all real $x$,
$f(x^2 + 1) = x^4 + 5x^2 + 3$.
For all real $x, f(x^2-1)$ is


 $\textbf{(A)}\ x^4+5x^2+1\qquad \textbf{(B)}\ x^4+x^2-3\qquad \textbf{(C)}\ x^4-5x^2+1\qquad \textbf{(D)}\ x^4+x^2+3\qquad \textbf{(E)}\ \text{none of these}$


 
 \textbf{Answer: }B



\textbf{Solution}

Let $y = x^2 + 1$. Then $x^2 = y - 1$, so we can write the given equation as 
 \begin{align*}f(y) &= x^4 + 5x^2 + 3 \\ &= (x^2)^2 + 5x^2 + 3 \\ &= (y - 1)^2 + 5(y - 1) + 3 \\ &= y^2 - 2y + 1 + 5y - 5 + 3 \\ &= y^2 + 3y - 1.\end{align*}
Then substituting $x^2 - 1$ for $y$, we get
 \begin{align*}f(x^2 - 1) &= (x^2 - 1)^2 + 3(x^2 - 1) - 1 \\ &= x^4 - 2x^2 + 1 + 3x^2 - 3 - 1 \\ &= x^4 + x^2 - 3.\end{align*}
The answer is therefore $\boxed{\textbf{(B)}}$.

 



\textbf{Solution 2}

Let $y=x^2.$ We have that 

 
 \begin{align*}f(x^2 + 1) &= x^4 + 5x^2 + 3 \\ &= y^2 + 5y + 3 \\ &= y^2 + 5y + 4 - 1 \\ &= (y+1)(y+4) -1 \\ &= (x^2 + 1)(x^2 + 4) - 1\end{align*}


 Thus, we have $f(n) = n(n+3)-1.$

 If we plug in $n = x^2-1$ we have

 
 \begin{align*}f(x^2 - 1) &= (x^2 -1)(x^2 + 2)-1 \\ &= (x^4 + x^2 - 2) - 1 \\ &= \boxed{x^4 + x^2 -3} \end{align*}


 


	\item (\textit{AHSME 1983}) Point $D$ is on side $CB$ of triangle $ABC$. If 
$\angle{CAD} = \angle{DAB} = 60^\circ\mbox{, }AC = 3\mbox{ and }AB = 6$, 
then the length of $AD$ is 


 $\textbf{(A)} \ 2 \qquad  \textbf{(B)} \ 2.5 \qquad  \textbf{(C)} \ 3 \qquad  \textbf{(D)} \ 3.5 \qquad  \textbf{(E)} \ 4$


 
 \textbf{Answer: }A



\textbf{Solution 1 (Law of Cosines)}

Let $AD = y$. Since $AD$ bisects $\angle{BAC}$, the Angle Bisector Theorem gives $\frac{DB}{CD} = \frac{AB}{AC} = 2$, so let $CD = x$ and $DB = 2x$. Applying the Law of Cosines to $\triangle CAD$ gives $x^2 = 3^2 + y^2 - 3y$, and to $\triangle DAB$ gives $(2x)^2 = 6^2 + y^2 - 6y$. Subtracting $4$ times the first equation from the second equation therefore yields $0 = 6y - 3y^2 \Rightarrow y(y-2) = 0$, so $y$ is $0$ or $2$. But since $y \neq 0$ ($y$ is the length of a side of a triangle), $y$ must be $2$, so the answer is $\boxed{\textbf{(A)} \ 2}$.

 



\textbf{Solution 2 (Similar Triangles)}

Let $E$ be a point on line $AC$ such that $AE = 6$ and $A$ is between $E$ and $C$.  Then $\triangle ABE$ must be equilateral, $BE = 6$, and $\angle BEC = 60^\circ$.  So $BE \parallel DA$ and $\triangle CDA \sim \triangle CBE$.  So $\frac{AD}{6} = \frac{3}{9}$, so $AD = \boxed{(\mathbf{A})\ 2}$.

 -j314andrews

 


	\item (\textit{AHSME 1983}) If $\tan{\alpha}$ and $\tan{\beta}$ are the roots of $x^2 - px + q = 0$, and $\cot{\alpha}$ and $\cot{\beta}$
are the roots of $x^2 - rx + s = 0$, then $rs$ is necessarily 


 $\textbf{(A)} \ pq \qquad  \textbf{(B)} \ \frac{1}{pq} \qquad  \textbf{(C)} \ \frac{p}{q^2} \qquad  \textbf{(D)}\ \frac{q}{p^2}\qquad \textbf{(E)}\ \frac{p}{q}$


 
 \textbf{Answer: }C



\textbf{Solution}

By Vieta's Formulae, we have $\tan(\alpha)\tan(\beta)=q$ and $\cot(\alpha)\cot(\beta)=s$. Recalling that $\cot\theta=\frac{1}{\tan\theta}$, we have $\frac{1}{\tan(\alpha)\tan(\beta)}=\frac{1}{q}=s$. 

 Also by Vieta's Formulae, we have $\tan(\alpha)+\tan(\beta)=p$ and $\cot(\alpha)+\cot(\beta)=r$, and again using $\cot\theta=\frac{1}{\tan\theta}$, we have $\tan(\alpha)+\tan(\beta)=r(\tan(\alpha)\tan(\beta))$. Using $\tan(\alpha)\tan(\beta)=q$ and $\tan(\alpha)+\tan(\beta)=p$, we therefore deduce that $r=\frac{p}{q}$, which yields $rs = \frac{1}{q}\cdot\frac{p}{q}=\frac{p}{q^2}$.

 Thus, the answer is $\boxed{\textbf{(C)} \ \frac{p}{q^2}}$.

 


	\item (\textit{AHSME 1983}) Find the smallest positive number from the numbers below.


 $\textbf{(A)} \ 10-3\sqrt{11} \qquad  \textbf{(B)} \ 3\sqrt{11}-10 \qquad  \textbf{(C)}\ 18-5\sqrt{13}\qquad \textbf{(D)}\ 51-10\sqrt{26}\qquad \textbf{(E)}\ 10\sqrt{26}-51$


 
 \textbf{Answer: }D



\textbf{Solution 1}

Notice that $3\sqrt{11} - 10 = \sqrt{99} - \sqrt{100} < 0$, $18 - 5\sqrt{13} = \sqrt{324} - \sqrt{325} < 0$, and similarly, $10\sqrt{26} - 51 = \sqrt{2600} - \sqrt{2601} < 0$, so these numbers can be excluded immediately as they are negative and we seek the smallest \textbf{positive} number. The remaining choices are $10 - 3\sqrt{11} = \sqrt{100} - \sqrt{99}$ and $51 - 10\sqrt{26} = \sqrt{2601} - \sqrt{2600}$. Observe that $\sqrt{100} - \sqrt{99} = \frac{1}{\sqrt{100} + \sqrt{99}} \approx \frac{1}{10+10} = \frac{1}{20}$, and $\sqrt{2601} - \sqrt{2600} = \frac{1}{\sqrt{2601} + \sqrt{2600}} \approx \frac{1}{51+51} = \frac{1}{102} < \frac{1}{20}$, so the smallest positive number in the list is $\boxed{\textbf{(D)} \ 51 - 10 \sqrt{26}}$.

 



\textbf{Solution 2}

After reaching the two possibilities $\sqrt{100}-\sqrt{99}$ and $\sqrt{2601}-\sqrt{2600}$ as in the first solution, note that these can be written as $\sqrt{x}-\sqrt{x-1}$ for $x=100$ and $x=2601$ respectively. Now, $\frac{\mathrm{d}}{\mathrm{d}x}\left(\sqrt{x}-\sqrt{x-1}\right) = \frac{1}{2\sqrt{x}} - \frac{1}{2\sqrt{x-1}} = \frac{\sqrt{x-1} - \sqrt{x}}{2\sqrt{x}\sqrt{x-1}} < 0$ (as the numerator is always negative and the denominator is always positive for $x \geq 1$). Therefore, with its derivative being negative, $\sqrt{x}-\sqrt{x-1}$ is a strictly decreasing function over its domain, so the smaller value will be the one with the larger value of $x$, namely $x = 2601$. Therefore the answer is $\boxed{\textbf{(D)} \ 51 - 10 \sqrt{26}}$ as before.

 


	\item (\textit{AHSME 1983}) Consider the two functions $f(x) = x^2+2bx+1$ and $g(x) = 2a(x+b)$, where the variable $x$ and the constants $a$ and $b$ are real numbers. 
Each such pair of constants $a$ and $b$ may be considered as a point $(a,b)$ in an $ab$-plane. 
Let $S$ be the set of such points $(a,b)$ for which the graphs of $y = f(x)$ and $y = g(x)$ do \textbf{not} intersect (in the $xy$-plane). The area of $S$ is 


 $\textbf{(A)} \ 1 \qquad  \textbf{(B)} \ \pi \qquad  \textbf{(C)} \ 4 \qquad  \textbf{(D)} \ 4 \pi \qquad  \textbf{(E)} \ \text{infinite}$


 
 \textbf{Answer: }B



\textbf{Solution}

We must describe geometrically those $(a,b)$ for which the equation $x^2+2bx+1=2a(x+b)$, i.e. $x^2+2(b-a)x+(1-2ab)=0$, has no solutions (equivalent to the graphs not intersecting). By considering the discriminant of this quadratic equation, there are no solutions if and only if $\left(2(b-a)\right)^2 - 4(1)(1-2ab) < 0 \Rightarrow (b-a)^2 < 1 - 2ab \Rightarrow a^2 - 2ab + b^2 < 1 - 2ab \Rightarrow a^2 + b^2 < 1$. Thus $S$ is the unit circle (without its boundary, due to the inequality sign being $<$ rather than $\leq$, but this makes no difference to the area), whose area is $\pi (1^2) = \pi$, so the answer is $\boxed{\textbf{(B)} \ \pi}$.

 


	\item (\textit{AHSME 1983}) In the adjoining figure the five circles are tangent to one another consecutively and to the lines 
$L_1$ and $L_2$. 
If the radius of the largest circle is $18$ and that of the smallest one is $8$, then the radius of the middle circle is


 
                    \begin{figure}[H]
                        centering
                        \includegraphics[width=0.5\textwidth]{Problem/AHSME_I_1983_23_Problem_Figure_1.png}
                    \end{figure}
                    


 $\textbf{(A)} \ 12 \qquad  \textbf{(B)} \ 12.5 \qquad  \textbf{(C)} \ 13 \qquad  \textbf{(D)} \ 13.5 \qquad  \textbf{(E)} \ 14$


 
 \textbf{Answer: }A



\textbf{Solution}


                    \begin{figure}[H]
                        centering
                        \includegraphics[width=0.5\textwidth]{Solution/AHSME_I_1983_23_Solution_Figure_1.png}
                    \end{figure}
                    

 Consider three consecutive circles, as shown in the diagram above; observe that their centres $P$, $Q$, and $R$ are collinear by symmetry. Let $A$, $B$, and $C$ be the points of tangency, and let $PS$ and $QT$ be segments parallel to the upper tangent (i.e. $L_1$), as also shown. Since $PQ$ is parallel to $QR$ (the three points are collinear), $PS$ is parallel to $QT$ (as both are parallel to $L_1$), and $SQ$ is parallel to $TR$ (as both are perpendicular to $L_1$, due to the tangent being perpendicular to the radius), we have $\triangle PQS \sim \triangle QRT$.

 Now, if we let $x, y$, and $z$ be the radii of the three circles (from smallest to largest), then $QS = y - x$ and $RT = z - y$. Thus, from the similarity that we just proved, $\frac{QS}{PQ} = \frac{RT}{QR} \Rightarrow \frac{y-x}{x+y} = \frac{z-y}{y+z}$ (where e.g. $PQ = x + y$ because of collinearity). This equation reduces to $y^2 = zx$, i.e. $\frac{y}{x} = \frac{z}{y}$, so the ratio of consecutive radii is constant, forming a geometric sequence. In this case, as the first radius is $8$ and, four radii later, the radius is $18$, this constant ratio is $\sqrt[\leftroot{-2}\uproot{2}{4}]{\frac{18}{8}}$. Therefore the middle radius is $8 \cdot {\left(\sqrt[\leftroot{-2}\uproot{2}{4}]{\frac{18}{8}}\right)}^{2} = 8 \sqrt{\frac{18}{8}} = \sqrt{18 \cdot 8} = \sqrt{144} = 12$, which is choice $\boxed{\textbf{(A)}}$.

 


	\item (\textit{AHSME 1983}) How many non-congruent right triangles are there such that the perimeter in $\text{cm}$ and the area in $\text{cm}^2$ are numerically equal? 


 $\textbf{(A)} \ \text{none} \qquad  \textbf{(B)} \ 1 \qquad  \textbf{(C)} \ 2 \qquad  \textbf{(D)} \ 4 \qquad  \textbf{(E)} \ \text{infinitely many}$


 
 \textbf{Answer: }E



\textbf{Solution 1}

Let the triangle have legs of length $a$ and $b$, so by the Pythagorean Theorem, the hypotenuse has length $\sqrt{a^2+b^2}$. Therefore we require 
 \begin{align*} &a + b + \sqrt{a^2+b^2} = \frac{1}{2} ab \\ \Rightarrow \quad &2 \sqrt{a^2+b^2} = ab - 2a - 2b \\ \Rightarrow \quad &4a^2 + 4b^2 = a^{2}b^{2} + 4a^2 + 4b^2 - 4a^{2}b - 4ab^{2} + 8ab \\ \Rightarrow \quad &4a^{2}b + 4ab^{2} = a^{2}b^{2} + 8ab. \end{align*}
 Now, as $a$ and $b$ are side lengths of a triangle, they must both be non-zero, so we can safely divide by $ab$ to give $4a + 4b = ab + 8 \Rightarrow b(a-4) = 4a-8 \Rightarrow b = \frac{4a-8}{a-4}$, so for any value of $a$ other than $4$, we can generate a valid corresponding value of $b$. 

 Notice also that each of these values of $a$ will give a \textbf{unique} corresponding value of $b$, since $\frac{4a-8}{a-4} = 4 + \frac{8}{a-4}$, and by considering the graph of $y = 4 + \frac{8}{x-4}$, it is clear that any horizontal line will intersect it at most once. Thus there are infinitely many valid solutions (one for every value of $a$ except $4$), so the answer is $\boxed{\textbf{(E)} \ \text{infinitely many}}$.

 



\textbf{Solution 2}

We use the formula $A = rs$, where $A$ is the area, $r$ is the inradius, and $s$ is the semiperimeter of a triangle. In this case, we have $A = 2s$, so as $A$ and $s$ are nonzero (the triangle must not be degenerate), we must have $r = 2$. Now simply observe that there are clearly many right triangles that can be drawn with inradius $2$ (try drawing a circle of radius $2$ tangent to two lines that form a right angle and then seeing how many lines tangent to the circle you can draw to form the third side of the triangle), so the answer must be $\boxed{\textbf{(E)} \ \text{infinitely many}}$ as before.

 


	\item (\textit{AHSME 1983}) If $60^a = 3$ and $60^b = 5$, then $12^{(1-a-b)/\left(2\left(1-b\right)\right)}$ is


 $\textbf{(A)} \ \sqrt{3} \qquad  \textbf{(B)} \ 2 \qquad  \textbf{(C)} \ \sqrt{5} \qquad  \textbf{(D)} \ 3 \qquad  \textbf{(E)} \ \sqrt{12}$


 
 \textbf{Answer: }B



\textbf{Solution 1}

We have that $12=\frac{60}{5}$. We can substitute our value for 5, to get
 \[12=\frac{60}{60^b}=60\cdot 60^{-b}=60^{1-b}.\]Hence
 \[12^{(1-a-b)/\left(2\left(1-b\right)\right)}=60^{(1-b)(1-a-b)/\left(2\left(1-b\right)\right)}=60^{(1-a-b)/2}.\]Since $4=\frac{60}{3\cdot 5}$, we have
 \[4=\frac{60}{60^a 60^b}=60^{1-a-b}.\]Therefore, we have
 \[60^{(1-a-b)/2}=4^{1/2}=\boxed{\textbf{(B)}\ 2}\]

 



\textbf{Solution 2}

We have $60^a = 3$ and $60^b = 5$. We can say that $a = \log_{60} 3$ and $b = \log_{60} 5$. 

 
 \[12^{(1-a-b)/[2(1-b)]} = 12^{(1-(a+b))/[2(1-b)]}\]

 We can evaluate (a+b) by the Addition Identity for Logarithms, $(a+b) = \log_{60} 15$. Also, $1 = \log_{60} 60$. 

 
 \[(1-(a+b) = \log_{60} 60 - \log_{60} 15 = \log_{60} 4\]

 Now we have to evaluate [2(1-b)], the denominator of the fractional exponent. Once again, we can say $1 = \log_{60} 60$

 
 \[2(1-b) = 2(\log_{60} 12)\]

 
 \[12^{(\log_{60} 4)/[2(\log_{60} 12]}  = 12^{\frac{1}{2} \cdot \log_{12} 4} = 4^{1/2} = 2\]

 


	\item (\textit{AHSME 1983}) The probability that event $A$ occurs is $\frac{3}{4}$; the probability that event $B$ occurs is $\frac{2}{3}$. 
Let $p$ be the probability that both $A$ and $B$ occur. The smallest interval necessarily containing $p$ is the interval


 $\textbf{(A)}\ \Big[\frac{1}{12},\frac{1}{2}\Big]\qquad \textbf{(B)}\ \Big[\frac{5}{12},\frac{1}{2}\Big]\qquad \textbf{(C)}\ \Big[\frac{1}{2},\frac{2}{3}\Big]\qquad \textbf{(D)}\ \Big[\frac{5}{12},\frac{2}{3}\Big]\qquad \textbf{(E)}\ \Big[\frac{1}{12},\frac{2}{3}\Big]$


 
 \textbf{Answer: }D



\textbf{Solution}

Firstly note that $p \leq \frac{3}{4}$ and $p \leq \frac{2}{3}$, as clearly the probability that both $A$ and $B$ occur cannot be more than the probability that $A$ or $B$ alone occurs. The more restrictive condition is $p \leq \frac{2}{3}$, since $\frac{2}{3} < \frac{3}{4}$. 

 Furthermore, by the \href{https://artofproblemsolving.com/wiki/index.php?title=Inclusion-Exclusion\_Principle}{Inclusion-Exclusion Principle}, we also have 
 \[\text{P}(A' \wedge B') = 1 - \text{P}(A) - \text{P}(B) + \text{P}(A \wedge B) = 1 - \frac{3}{4} - \frac{2}{3} + p = p - \frac{5}{12},\] and as a probability must be non-negative, $p - \frac{5}{12} \geq 0$, so $p \geq \frac{5}{12}$. Therefore, combining our inequalities gives $\frac{5}{12} \leq p \leq \frac{2}{3}$, or $\boxed{\textbf{(D)} \Big[\frac{5}{12},\frac{2}{3}\Big]}$.

 


	\item (\textit{AHSME 1983}) A large sphere is on a horizontal field on a sunny day. At a certain time the shadow of the sphere reaches out a distance 
of $10$ m from the point where the sphere touches the ground. At the same instant a meter stick 
(held vertically with one end on the ground) casts a shadow of length $2$ m. What is the radius of the sphere in meters? 
(Assume the sun's rays are parallel and the meter stick is a line segment.)


 $\textbf{(A)}\ \frac{5}{2}\qquad \textbf{(B)}\ 9 - 4\sqrt{5}\qquad \textbf{(C)}\ 8\sqrt{10} - 23\qquad \textbf{(D)}\ 6-\sqrt{15}\qquad \textbf{(E)}\ 10\sqrt{5}-20$


 
 \textbf{Answer: }E



\textbf{Solution}


                    \begin{figure}[H]
                        centering
                        \includegraphics[width=0.5\textwidth]{Solution/AHSME_I_1983_27_Solution_Figure_1.png}
                    \end{figure}
                    

 Let $O$ be the center of the sphere, $T$ be the point where the sphere touches the ground, $S$ be the tip of the sphere's shadow, $R$ be the point where sun rays that reach the ground near $S$ would approach the sphere, and $\theta = \angle RST$ be the angle between the sun's rays and the ground. 

 Then $RS$ and $RT$ must both be tangent to the sphere, and therefore $RS = ST = 10$ and $OS$ bisects $\angle RST$.  So $\angle OST = \frac{\theta}{2}$.

 Since a $1$ meter stick casts a $2$ meter shadow, $\tan \theta = \frac{1}{2}$.  So $\cos \theta = \frac{1}{\sec \theta} = \frac{1}{\sqrt{1+\tan^2 \theta}} = \frac{1}{\sqrt{\frac{5}{4}}} = \frac{2}{\sqrt{5}}$, and $\sin \theta = \sqrt{1 - \cos^2 \theta} = \sqrt{\frac{1}{5}} = \frac{1}{\sqrt{5}}$.  So $\tan \left(\frac{\theta}{2}\right) = \frac{1-\cos \theta}{\sin \theta} = \frac{1 - \frac{2}{\sqrt{5}}}{\frac{1}{\sqrt{5}}} = \sqrt{5}-2$.

 Let $r = OT = OR$ be the radius of the sphere.  Then $\tan\left(\frac{\theta}{2}\right) = \frac{r}{10}$, so $r = 10 \tan\left(\frac{\theta}{2}\right) = 10(\sqrt{5}-2) = \boxed{(\mathbf{E})\ 10\sqrt{5} - 20}$.

 


	\item (\textit{AHSME 1983}) Triangle $\triangle ABC$ in the figure has area $10$. Points $D, E$ and $F$, all distinct from $A, B$ and $C$, 
are on sides $AB, BC$ and $CA$ respectively, and $AD = 2, DB = 3$. If triangle $\triangle ABE$ and quadrilateral $DBEF$ 
have equal areas, then that area is


 
                    \begin{figure}[H]
                        centering
                        \includegraphics[width=0.5\textwidth]{Problem/AHSME_I_1983_28_Problem_Figure_1.png}
                    \end{figure}
                    


 $\textbf{(A)}\ 4\qquad \textbf{(B)}\ 5\qquad \textbf{(C)}\ 6\qquad \textbf{(D)}\ \frac{5}{3}\sqrt{10}\qquad \textbf{(E)}\ \text{not uniquely determined}$


 
 \textbf{Answer: }C



\textbf{Solution}

Let $G$ be the intersection point of $AE$ and $DF$. Since $[DBEF] = [ABE]$, we have $[DBEG] + [EFG] = [DBEG] + [ADG]$, i.e. $[EFG] = [ADG]$. It therefore follows that $[ADG] + [AGF] = [EFG] + [AGF]$, so $[ADF] = [AFE]$. Now, taking $AF$ as the base of both $\triangle ADF$ and $\triangle AFE$, and using the fact that triangles with the same base and same perpendicular height have the same area, we deduce that the perpendicular distance from $D$ to $AF$ is the same as the perpendicular distance from $E$ to $AF$. This in turn implies that $AF \parallel DE$, and so as $A$, $F$, and $C$ are collinear, $AC \parallel DE$. Thus $\triangle DBE \sim \triangle ABC$, so $\frac{BE}{BC} = \frac{BD}{BA} = \frac{3}{5}$. Since $\triangle ABE$ and $\triangle ABC$ have the same perpendicular height (taking $AB$ as the base), $[ABE] = \frac{3}{5} \cdot [ABC]=\frac{3}{5}\cdot 10 = 6$, and hence the answer is $\boxed{\textbf{(C)}}$.

 


	\item (\textit{AHSME 1983}) A point $P$ lies in the same plane as a given square of side $1$. Let the vertices of the square, 
taken counterclockwise, be $A, B, C$ and $D$. Also, let the distances from $P$ to $A, B$ and $C$, respectively, be $u, v$ and $w$. 
What is the greatest distance that $P$ can be from $D$ if $u^2 + v^2 = w^2$?


 $\textbf{(A)}\ 1 + \sqrt{2}\qquad \textbf{(B)}\ 2\sqrt{2}\qquad \textbf{(C)}\ 2 + \sqrt{2}\qquad \textbf{(D)}\ 3\sqrt{2}\qquad \textbf{(E)}\ 3+\sqrt{2}$


 
 \textbf{Answer: }C



\textbf{Solution}

Place the square in the $xy$-plane with $A$ as the origin, so that $B=(1,0), C=(1,1),$ and $D=(0,1).$ We are given that $PA^2+PB^2=PC^2,$ so 

 
 \begin{align*}&(x^2+y^2)+((x-1)^2+y^2)=(x-1)^2+(y-1)^2\\ \Rightarrow \quad &2x^2+2y^2-2x+1=x^2+y^2-2x-2y+2\\ \Rightarrow \quad &x^2+y^2=-2y+1\\ \Rightarrow \quad &x^2+y^2+2y-1=0\\ \Rightarrow \quad &x^2+(y+1)^2=2.\end{align*}


 Thus we see that $P$ lies on a circle centered at $(0,-1)$ with radius $\sqrt{2}.$ The farthest point from $D$ on this circle is at the bottom of the circle, at $\left(0, -1-\sqrt{2}\right),$ in which case $PD$ is $1 - \left(-1 - \sqrt{2}\right) = \boxed{\textbf{(C)}\ 2 + \sqrt{2}}.$

 


	\item (\textit{AHSME 1983}) Distinct points $A$ and $B$ are on a semicircle with diameter $MN$ and center $C$. 
The point $P$ is on $CN$ and $\angle CAP = \angle CBP = 10^{\circ}$. If $\stackrel{\frown}{MA} = 40^{\circ}$, then $\stackrel{\frown}{BN}$ equals


 
                    \begin{figure}[H]
                        centering
                        \includegraphics[width=0.5\textwidth]{Problem/AHSME_I_1983_30_Problem_Figure_1.png}
                    \end{figure}
                    


 $\textbf{(A)}\ 10^{\circ}\qquad \textbf{(B)}\ 15^{\circ}\qquad \textbf{(C)}\ 20^{\circ}\qquad \textbf{(D)}\ 25^{\circ}\qquad \textbf{(E)}\ 30^{\circ}$


 
 \textbf{Answer: }C



\textbf{Solution 1}

Since $\angle CAP = \angle CBP = 10^\circ$, quadrilateral $ABPC$ is cyclic (as shown below) by the converse of the theorem "angles inscribed in the same arc are equal".

 
                    \begin{figure}[H]
                        centering
                        \includegraphics[width=0.5\textwidth]{Solution/AHSME_I_1983_30_Solution_Figure_1.png}
                    \end{figure}
                    

 Since $\angle ACM = 40^\circ$, $\angle ACP = 140^\circ$, so, using the fact that opposite angles in a cyclic quadrilateral sum to $180^{\circ}$, we have $\angle ABP = 40^\circ$. Hence $\angle ABC = \angle ABP - \angle CBP = 40^ \circ - 10^\circ = 30^\circ$.

 Since $CA = CB$, triangle $ABC$ is isosceles, with $\angle BAC = \angle ABC = 30^\circ$. Now, $\angle BAP = \angle BAC - \angle CAP = 30^\circ - 10^\circ = 20^\circ$. Finally, again using the fact that angles inscribed in the same arc are equal, we have $\angle BCP = \angle BAP = \boxed{\textbf{(C)}\ 20^{\circ}}$.

 



\textbf{Solution 2}


                    \begin{figure}[H]
                        centering
                        \includegraphics[width=0.5\textwidth]{Solution/AHSME_I_1983_30_Solution_Figure_2.png}
                    \end{figure}
                    

 Complete the circle centered at $C$ and extend $AP$ and $AC$ to meet the circle at $Q$ and $R$, respectively.  So $\stackrel{\frown}{NR} = \stackrel{\frown}{MA} = 40^\circ$, and $\stackrel{\frown}{QR} = 2 \cdot 10^\circ = 20^\circ$.  Therefore, $\stackrel{\frown}{NQ} = 40^\circ - 20^\circ = 20^\circ$.  

 Since $CA = CQ$, $\triangle CAQ$ is isosceles and $\angle CQA = \angle CAQ = 10^\circ$.  Let $\theta = \angle CBQ = \angle CQB$.  Then $\angle PBQ = \angle PQB = \theta - 10^\circ$, so $\triangle PBQ$ is isosceles with $PB = PQ$.  So $CN$ must bisect $\angle BCQ$, and $\stackrel{\frown}{BN} = \stackrel{\frown}{NQ} = \boxed{(\mathbf{C})\ 20^{\circ}}$.

 -j314andrews

 


\end{enumerate}